\documentclass[11pt]{article}
\usepackage{amsmath, amssymb, amsthm}
\usepackage{geometry}
\usepackage{hyperref}
\usepackage{parskip}
\geometry{margin=1.2in}

\newcommand{\R}{\mathbb{R}}
\newcommand{\C}{\mathbb{C}}
\newcommand{\Fhat}{\hat{F}}
\newcommand{\psihat}{\hat{\psi}}
\newcommand{\imag}{\operatorname{Im}}
\newcommand{\real}{\operatorname{Re}}

\title{Notes on Kotlarski and Frequency Marching:\\
       Zero Self-Correction and Quadrature Noise Sensitivity}
\date{}

\begin{document}
\maketitle

\section{Setup and the Sign Ambiguity}

Let $f$ be a real-valued signal with Fourier transform $\Fhat(\omega) =
|\Fhat(\omega)|\,e^{i\varphi(\omega)}$.  Because $f$ is real,
$\Fhat(-\omega) = \overline{\Fhat(\omega)}$, so the full spectrum is
determined by the positive-frequency half.  In the MRA problem we
estimate the \emph{bispectrum}
\[
    B(i,j) \;\propto\; \Fhat(i)\,\overline{\Fhat(j)}\,\overline{\Fhat(i-j+1)},
\]
from which we wish to recover $\varphi$.  Taking the angle,
\[
    \Psi(i,j) := \angle B(i,j) = \varphi(i) - \varphi(j) - \varphi(i-j+1)
    \pmod{2\pi}.
\]

The fundamental difficulty is that where $\Fhat(\omega) \approx 0$,
the phase $\varphi(\omega)$ is numerically ambiguous: a small perturbation
can flip the sign of $\Fhat$, shifting $\varphi$ by $\pm\pi$.  Any
algorithm that passes through such a zero must decide how to handle
this $\pm\pi$ jump.

\section{Diagonal Kotlarski}

Kotlarski's approach recovers $\log\Fhat(\omega)$ by integrating the
log-derivative of $\psihat$ along the diagonal:
\[
    \log\Fhat(\omega_k) = \int_0^{\omega_k}
    \frac{\partial}{\partial\omega_1}\log\psihat(\omega_1,\omega_1)\,d\omega_1
    + \text{const}.
\]
This is a \textbf{single cumulative path integral}.  If $\Fhat$ has a
zero at $\omega_z < \omega_k$, the imaginary part of the integrand
contains a Dirac-like $\pm\pi$ spike at $\omega_z$.  After integration,
\emph{every} frequency $\omega_k > \omega_z$ carries the accumulated
$\pm\pi$ offset.  There is only one path, so there is no other estimate
to compare against.  The error persists for all higher frequencies until
it is explicitly detected and corrected---this is zero-flipping.

\section{Frequency Marching}

Frequency marching (FM) builds up $\varphi(i)$ for $i = 3, 4, \ldots$
using the bispectral phase identity.  At step $i$, it forms multiple
estimates:
\[
    \widehat\varphi^{(j)}(i) = \Psi(i, j) + \psi(j) + \psi(i-j+1),
    \qquad j = 2, \ldots, \lceil i/2 \rceil,
\]
where $\psi(k)$ denotes the already-computed estimate of $\varphi(k)$.
Writing $\varepsilon(k) = \psi(k) - \varphi(k)$ for the error at
frequency $k$, each estimate satisfies
\[
    \widehat\varphi^{(j)}(i)
    = \varphi(i) + \varepsilon(j) + \varepsilon(i-j+1).
\]

The final estimate is the \textbf{circular mean}:
\[
    \psi(i) = \angle\!\left(\sum_{j=2}^{\lceil i/2\rceil}
    e^{i\,\widehat\varphi^{(j)}(i)}\right).
\]

\subsection*{Why this self-corrects at zeros}

Suppose $\Fhat$ has a single zero at frequency $z$, so
$\varepsilon(z) \approx \pm\pi$ while $\varepsilon(k) \approx 0$ for
$k \neq z$.  For a given target frequency $i$, path $j$ carries a
$\pm\pi$ error if and only if $j = z$ or $j = i-z+1$.  At most
\textbf{two} paths are corrupted out of roughly $i/2$ total paths.

On the unit circle:
\begin{itemize}
    \item Each \emph{good} path contributes $e^{i\varphi(i)}$.
    \item Each \emph{bad} path contributes $e^{i(\varphi(i)\pm\pi)}
          = -e^{i\varphi(i)}$.
\end{itemize}
The vector sum is approximately
\[
    \left(\frac{i}{2} - 2\right)e^{i\varphi(i)}
    - 2\,e^{i\varphi(i)}
    = \left(\frac{i}{2} - 4\right)e^{i\varphi(i)},
\]
which points in the direction $e^{i\varphi(i)}$ for $i > 8$.  The
circular mean therefore recovers $\varphi(i)$ correctly despite the
corrupted paths.  As $i$ grows, the fraction of bad paths shrinks
as $O(1/i)$, so the correction becomes increasingly reliable.

Moreover, because $\psi(i)$ is corrected at each step, the error does
\emph{not} propagate forward: future steps receive a clean $\psi(i)$
as input.  This is the key structural difference from Kotlarski.

\subsection*{Why single-path FM fails}

If the inner loop is restricted to $j = 2$ only, the estimate at step
$i$ is:
\[
    \psi(i) = \Psi(i, 2) + \psi(2) + \psi(i-1).
\]
Unrolling the recursion on $\psi(i-1), \psi(i-2), \ldots$:
\[
    \psi(i) = \sum_{k=3}^{i} \Psi(k, 2) + (i-2)\,\psi(2) + \psi(2).
\]
This is a \textbf{cumulative sum along a fixed path} through the
bispectrum---structurally identical to diagonal Kotlarski.  A zero at
any frequency $z < i$ contributes a $\pm\pi$ error to every $\psi(k)$
for $k > z$, with no averaging to correct it.

The multi-path circular average breaks this structure.  At each step
$i$, the estimate is formed fresh from paths of all lengths $j =
2,\ldots,\lceil i/2\rceil$, rather than being chained through the
previous step.  The self-correcting property comes entirely from having
enough independent paths at each step to outvote the corrupted ones.

\section{Multipath Kotlarski}

In the multipath Kotlarski algorithm, each estimate of $g(k)$ (the
log-amplitude or phase at positive frequency $k$) is:
\[
    \widehat g^{(c)}(k) = g(c) + I_c(k) + g(k-c),
\]
where $I_c(k)$ is a \emph{path integral} of the bispectral
log-derivative along column $c$ of the bispectrum from row $c$ to row
$k$.  This resembles FM in that it averages over multiple paths, and
in the phase-only (``mixed'') mode one can replace the arithmetic mean
with a circular mean to gain some robustness.

However, there is an important difference.  In FM, each estimate
$\widehat\varphi^{(j)}(i)$ depends on two previously computed scalars
$\psi(j)$ and $\psi(i-j+1)$; the error structure is clean.  In
multipath Kotlarski, $I_c(k)$ itself is an integral that may
\emph{internally cross a zero} of $\Fhat$ between $\omega_c$ and
$\omega_k$, picking up a $\pm\pi$ contribution that is not simply
inherited from a prior step but arises fresh from within the integral.
The set of corrupted paths is therefore less predictable than in FM.

Circular averaging helps when the majority of paths are good, and for
functions with isolated zeros it can substantially reduce the need for
explicit zero-flipping.  But it is not a complete substitute: FM's
self-correction works because its estimates are structurally
non-accumulated (one step of bispectral phase per estimate), while
Kotlarski's integration paths always carry some accumulated history.

\section{Numerical Integration Order and Noise Sensitivity}

A subtlety that arises in practice is that higher-order quadrature rules
do not always produce better results when the integrand is noisy.  In the
Kotlarski integral
\[
    g(\omega_k) = \int_0^{\omega_k}
    \frac{\partial_{\omega_1}\psihat(\omega_1,\omega_1)}{\psihat(\omega_1,\omega_1)}\,d\omega_1,
\]
the integrand is estimated from finite data and therefore carries
pointwise noise.  The choice of quadrature rule determines how that noise
is weighted.

\subsection*{Riemann sum}

The simplest approximation evaluates the integrand at one point per
interval and weights each point equally:
\[
    g(\omega_k) \approx \Delta\omega \sum_{j=1}^{k} h(\omega_j),
    \qquad \text{weight on each } h(\omega_j) = 1.
\]

\subsection*{Simpson's rule}

Simpson's rule achieves $O((\Delta\omega)^4)$ accuracy for smooth
integrands by fitting local quadratics.  Over each pair of intervals it
uses the three-point formula
\[
    \int_{\omega_{j-1}}^{\omega_{j+1}} h\,d\omega
    \approx \frac{\Delta\omega}{3}\bigl(h(\omega_{j-1})
    + 4\,h(\omega_j) + h(\omega_{j+1})\bigr).
\]
The interior point $\omega_j$ receives a weight of $4/3\,\Delta\omega$,
while the endpoint weights are $1/3\,\Delta\omega$.  Interior points are
therefore weighted \textbf{four times more heavily} than they would be
under the Riemann sum.

\subsection*{Why Simpson's rule can be worse under noise}

The increased weight on interior points is precisely what gives Simpson's
rule its superior accuracy for smooth functions---but it backfires when
the integrand contains noise.  If the noise at frequency $\omega_j$ has
variance $\sigma^2$, Simpson's rule contributes variance
$(4\Delta\omega/3)^2\sigma^2 \approx 16/9\,(\Delta\omega)^2\sigma^2$ from
that point, compared to $(\Delta\omega)^2\sigma^2$ under the Riemann sum.
Every interior point is amplified by a factor of $16/9 \approx 1.78$ in
variance relative to the Riemann sum.

Because the Kotlarski integral is cumulative---each $g(\omega_k)$ depends
on all quadrature evaluations from $0$ to $\omega_k$---these amplified
noise contributions accumulate along the integration path.  At low noise
levels the smoothness gain of Simpson's rule dominates.  At higher noise
levels the amplification dominates, and the Riemann sum (or trapezoidal
rule, which also uses weight 1 at interior points) produces smaller
accumulated error.

The trapezoidal rule sits between the two: it also uses weight 1 at
interior points and achieves $O((\Delta\omega)^2)$ accuracy for smooth
functions, making it generally the recommended default for noisy
integrands.  The practical takeaway is that the optimal quadrature order
depends on the signal-to-noise ratio of the integrand: higher-order rules
require the integrand to be well-resolved relative to the noise floor.

\subsection*{Formal variance comparison: trapezoidal vs.\ Simpson's rule}

Suppose the integrand is evaluated at $N+1$ equally-spaced points
$\omega_0 < \omega_1 < \cdots < \omega_N$ with spacing $h = \Delta\omega$,
and that each evaluation carries independent noise with variance $\sigma^2$.
The variance of any linear quadrature estimate $\hat I = \sum_j w_j\,h(\omega_j)$
is
\[
    \operatorname{Var}(\hat I) = \sigma^2 \sum_{j=0}^{N} w_j^2.
\]
We compare the two rules by computing $\sum_j w_j^2$ in each case.

\paragraph{Trapezoidal rule.}
The weights are $w_0 = w_N = h/2$ and $w_j = h$ for $j = 1,\ldots,N-1$.
\[
    \sum_{j=0}^N w_j^2
    = 2\!\left(\frac{h}{2}\right)^{\!2} + (N-1)\,h^2
    = h^2\!\left(N - \tfrac{1}{2}\right)
    \;\approx\; Nh^2.
\]

\paragraph{Composite Simpson's rule.}
Assuming $N$ is even, the weights are: $w_0 = w_N = h/3$; $w_j = 4h/3$
for $j$ odd ($N/2$ such points); and $w_j = 2h/3$ for $j$ even and interior
($N/2 - 1$ such points).
\begin{align*}
    \sum_{j=0}^N w_j^2
    &= 2\!\left(\frac{h}{3}\right)^{\!2}
       + \frac{N}{2}\!\left(\frac{4h}{3}\right)^{\!2}
       + \left(\frac{N}{2}-1\right)\!\left(\frac{2h}{3}\right)^{\!2} \\
    &= \frac{h^2}{9}\!\left[2 + 8N + 2N - 4\right]
     = \frac{h^2}{9}(10N - 2)
    \;\approx\; \frac{10}{9}\,Nh^2.
\end{align*}

\paragraph{Ratio.}
\[
    \frac{\operatorname{Var}_{\mathrm{Simpson}}}{\operatorname{Var}_{\mathrm{trap}}}
    \;\approx\; \frac{10/9\;Nh^2}{Nh^2} = \frac{10}{9} \approx 1.11.
\]
Composite Simpson's rule therefore has roughly \textbf{11\% higher noise
variance} than the trapezoidal rule, regardless of $N$ or $h$.  The
source is the large midpoint weights $4h/3$: each contributes
$(4h/3)^2 = 16h^2/9$ to the sum, versus $h^2$ per interior point under
the trapezoidal rule.

For the cumulative Kotlarski integral the penalty compounds: the estimate
at $\omega_k$ sums quadrature weights from $\omega_0$ to $\omega_k$, so
the 11\% excess variance at each step accumulates along the path.
Combined with the fact that MATLAB's built-in \texttt{cumtrapz} is
substantially faster than an interpreted Simpson's implementation
(which also involves a Vandermonde solve for even-length inputs), the
trapezoidal rule is the preferred quadrature method whenever the
integrand is noisy.

\section{Summary}

\begin{center}
\begin{tabular}{lll}
\hline
\textbf{Algorithm} & \textbf{Path structure} & \textbf{Zero handling} \\
\hline
Diagonal Kotlarski & Single cumulative integral & Must flip explicitly \\
Single-path FM ($j{=}2$) & Cumulative sum in disguise & Must flip explicitly \\
Multi-path FM & Fresh estimates + circular mean & Self-corrects (majority vote) \\
Multipath Kotlarski & Averaged integrals + circular mean & Partially self-corrects \\
\hline
\end{tabular}
\end{center}

The reason multi-path FM avoids explicit zero-flipping is not that it
avoids error propagation in general---FM still marches and can accumulate
errors---but that the \textbf{circular averaging at each step resets the
estimate} before it is used as input to the next step.  A single path
through the bispectrum, whether in Kotlarski or FM, gives no such reset
and is therefore vulnerable to uncorrected sign flips at every zero it
encounters.

\end{document}
